% >>> typage-template article@2026.10.1 variant=a4
\documentclass[11pt,a4paper]{article}
\usepackage{fontspec}
\usepackage[margin=2.5cm]{geometry}
\bibliographystyle{plain}
% <<< typage-template
\usepackage{amsmath,amssymb}
\usepackage{graphicx}
\usepackage{booktabs}
\usepackage[hidelinks]{hyperref}

\title{Article Title}
\author{First Author\thanks{Department, University, \texttt{first.author@example.com}}
  \and Second Author\thanks{Laboratory, Institute, \texttt{second.author@example.com}}}
\date{\today}

\begin{document}
\maketitle

\begin{abstract}
State the problem, the approach and the main result in a few sentences.
\end{abstract}

\section{Introduction}
Introduce the problem and why it matters, then list the contributions.
Typesetting follows the usual \LaTeX{} conventions~\cite{lamport1994latex}.

\section{Method}
Describe the method. Number the equations you refer to:
\begin{equation}
  \label{eq:objective}
  \min_{\theta} \; \frac{1}{n} \sum_{i=1}^{n} \ell\bigl(f_\theta(x_i), y_i\bigr) + \lambda \lVert \theta \rVert_2^2 .
\end{equation}
Equation~\eqref{eq:objective} is the training objective.

\section{Results}
Table~\ref{tab:results} summarizes the results.

\begin{table}[t]
  \centering
  \caption{Results on the example benchmark.}
  \label{tab:results}
  \begin{tabular}{lcc}
    \toprule
    Method   & Accuracy (\%) & Time (s) \\
    \midrule
    Baseline & 81.2          & 12.4     \\
    Ours     & 84.7          & 10.9     \\
    \bottomrule
  \end{tabular}
\end{table}

\section{Conclusion}
Conclude and outline future work.

\bibliography{refs}
\end{document}
